Discuss finite differences. What is the advantage of this method compared to other
methods, what is the disadvantage? Why is the central difference estimator pre-
ferred to the standard difference quotient? How should the discretization depend
on the number of samples? What is the best possible rate of the root mean-squared
error? How does this compare to standard MC?


Given is a parameter \( ϑ \).
The random variable \( Z(ϑ) \) depends on \( ϑ \) and can be simulated.
The Monte Carlo estimation of \( z(ϑ) = \mathbb{E}[Z(ϑ)] \) is tractable.
Assumption: \( ϑ \mapsto z(ϑ) \) is smooth, i. e., the function \( z \) is multiple times continuously differentiable.


If \(f(\vartheta)\) is smooth, we know that:
\[ \lim_{h \to 0} \frac{f(\vartheta + h) - f(\vartheta)}{h} = \lim_{h \to 0} \frac{f(\vartheta + \frac{h}{2}) - f(\vartheta - \frac{h}{2})}{h} \]

A Taylor expansion provides us with the following equations:

\( \Delta_f = \frac{f(\vartheta + h) - f(\vartheta)}{h} = f'(\vartheta) + \frac{h}{2} f''(\vartheta) + \mathcal{O}(h^2)\) 
which is the forward-difference estimator.

While
\(= \frac{f(\vartheta + \frac{h}{2}) - f(\vartheta - \frac{h}{2})}{h} = f'(\vartheta) + \frac{h^2}{24} f'''(\vartheta) + \mathcal{O}(h^4)\)
is the central difference estimator

The preferred random variable for the construction of the estimator of the derivative is
\( D(\vartheta) = \frac{Z(\vartheta + \frac{h}{2}) - Z(\vartheta - \frac{h}{2})}{h} \)

Sample \( Z(\vartheta + \frac{h_R}{2}), Z(\vartheta - \frac{h_R}{2}) \)
R-times independently for each replication.
Consider \( R \uparrow \infty \)  and a corresponding choice \( h_R \downarrow 0 \).
How should the function \( h_R \) be chosen to minimize the rate of the mean-square error?

\(\mathbb{E}[(\hat{Z}'(\vartheta) - Z'(\vartheta) )^2] \to 0 \)
when \( R \to \infty \) with \( \hat{Z}'(\vartheta) \) 
based on \( R \) independent replications of \( D(\vartheta) \)  using \( h_R \)? 

One can show that \( h_R = \mathcal{O}(R^{-1/6}) \) leads to the best possible rate \(  \mathcal{O}(R^{-1/3}) \)   for independent sampling.
But a better rate can be achieved by using common random numbers if \( Z(\vartheta) \) is smooth a.s.

Finite differences allow higher-order expansions that improve the optimal rates. But these methods implicitly involve FD approximations of
higher-order derivatives requiring the evaluation of z at multiple points. This increases computational effort and may also affect the
variance. Additionally the optimal choice of samples \( n \) vs. step size \( h \) is a careful balance,
since the first influences the variance and the later variance and bias (which is non-zero)
In practice, these methods are thus typically not used.

What is a kernel estimator of a function? How can you construct an estimator of
the derivative from this?

Estimation of \(z(\vartheta)\): more weight of samples close to \(\vartheta\) than further away, according to kernel \(k\)
Estimation of \(z'(\vartheta)\): exploit smoothness of \(k\)
k kernel: typically smooth and symmetric probability density,
e. g., normal density

Kernel estimator of  \(z(\vartheta)\): Assume \(\vartheta \in [0,1]\)

\[ \frac{1}{(R+1)h} \sum_{r=0}^{R} Z \left( \frac{r}{R} \right) k \left( \frac{\vartheta - r/R}{h} \right), \quad h > 0 \]

Observations:
For \(h\) smaller, the term \(Z \left( \frac{r}{R} \right)\) with \(\frac{r}{R}\) far away from \(\vartheta\) gets less weight. \(h\) is called bandwidth.
\[ \frac{1}{(R+1)h} \sum_{r=0}^{R} k \left( \frac{\vartheta - r/R}{h} \right) \]
is a Riemann-sum approximating \(\int k(y) dy = 1\) with a partition of size \(\frac{1}{Rh}\) and with \(R+1\) buckets.
Convergence requires \(R \uparrow \infty, h \neq 0\) and \(Rh \rightarrow \infty\). 

Exploiting the smoothness of k, differentiation with respect to \(\vartheta\) yields
\(D_{R}(\vartheta)=\frac{1}{(R+1)h^{2}}\sum_{r=0}^{R}Z \left( \frac{r}{R} \right) k^{\prime} \left( \frac{\vartheta - r/R}{h} \right). \)

Letting first tend \(R \uparrow \infty\) and then \(h \downarrow 0\), 
one observes that convergence to \(z'(\vartheta)\) occurs if \(k(x) = \mathcal{O}(|x|)\) as 
\(x \rightarrow \infty\).
The value of the bandwidth \(h > 0\) can be chosen as a function of \(R \uparrow \infty\), but as in the case of finite differences the rate is significantly slower than \(R^{-1/2}\). 


Samplewise / pathwise differentiation leads to infinitesimal perturbation analysis.
Which conditions are necessary for this to work? What is the rate of convergence
of the corresponding MC estimator?

Infinitesimal Perturbation Analysis (IPA)

Based on sample path differentiation
* For almost all \(\omega \in \Omega\), \(\vartheta \mapsto Z(\vartheta)(\omega)\) is a smooth function.
* Thus, \(D(\vartheta)(\omega) = \frac{d}{d\vartheta} Z(\vartheta)(\omega)\) can be computed pathwise.
The goal is to estimate \(z'(\vartheta)\) With \(z(\vartheta) = \mathbb{E}[Z(\vartheta)]\).
Estimator - 
\[
\hat{z}'_R(\vartheta) = \frac{D_1 + ... + D_R}{R} \text{ with } D_1, ..., D_R \sim D \text{ iid}
\]
The requirement for consistency (LLN) is:
\[\mathbb{E} \left[ \frac{d}{d\vartheta} Z(\vartheta) \right] = \mathbb{E} \left[ \frac{d}{d\vartheta} Z(\vartheta) \right] (*)\]
and the convergence rate by CLT:
\[\sqrt{R} (\hat{z}'_R(\vartheta) - z'(\vartheta)) \sim \mathcal{N}(0, \text{Var}(D)),\]
i.e., the estimator converges at the canonical rate \(R^{-1/2}\). 

If \( \vartheta \mapsto z(\vartheta) \) can be expressed on a common probability space and the function is smooth, often IPA can be
used.
In this case, also FD can be implemented with \( Z(ϑ + \frac{h}{2} ), Z (ϑ − \frac{h}{2}) \) sampled not independently, but on the
basis of common random numbers.
If IPA is admissible, then the preferred choice is IPA. If IPA is not admissible, common random numbers can
improve the performance of FD.

Discuss two simple examples involving two independent random variables \( Z_1, Z_2 \)
that admit respectively do not admit the application of infinitesimal perturbation
analysis.

This also answers the question:
Explain in detail how infinitesimal perturbation analysis can be used to compute
the sensitivity of a European-type derivative price with respect to the initial value
of the asset.

Imagine a network consisting of nodes \( n_1, n_2, n_3 \), where \( n_1 \) is the source
and \( n_3 \) the sink with connecting edges \( e_1 = (n_1, n_2), e_2 = (n_2, n_3)\). 
Let \( X_1, \varthetaX_2 \), where \( \vartheta \in [0, 1] \) be the random variables specifying the capacity of \( e_i \)
and the capacity of the whole network as \( \text{max}\{X_1, \vartheta X_2\} = Z(\theta)  \)
Then \( D(\vartheta) = 0 \) if \( X_1 > \vartheta X_2 \) and \( D(\vartheta = X_2 \) if \( X_1 < \vartheta X_2 \),
while it is undefined for \( \vartheta = \frac{X_1}{X_2} \).
So for this \( Z \) the IPA works.

If we assume that \( \vartheta X_2 > X_1  \) and we define \( Z(\vartheta = 1_{\vartheta X_2 > X_1} \)
then \( D(\vartheta) = 0 \) for all \( \vartheta \neq \frac{X_1}{X_2} \). If \( \mathbb{P}[X_1 = X_2] = 0 \), \( D(\vartheta) = 0 \) a.s.
But \( z(\vartheta) =  \mathbb{P}[\vartheta X_2 > X_1]\) typically possesses a derivative \( z'(\vartheta) \neq 0 \).
Thus the IPA method is not applicable here, since 
\( z'(\vartheta) = \frac{d}{d\vartheta}\mathbb{E}[Z(\vartheta)] \neq 0\)
and
\( \mathbb{E}[\frac{d}{d\vartheta}Z(\vartheta)] = D(\vartheta) = 0 \) a.s.
do not match.

Explain in detail how infinitesimal perturbation analysis can be used to compute
the sensitivity of a European-type derivative price with respect to the initial value
of the asset.

IPA works for the following setup, mainly because Lipschitz-continuity is satisfied

Underlying asset price process: 
\( (X(t))_{0 < t \leq T} = (X(t, x))_{0 < t \leq T} \) for \(X(0) = X(0, x) = x > 0 \)

Common model for any \( x > 0 \):

\[ dX(t) = a(X(t)) dt + b(X(t)) dB(t), \quad X(0) = x \]

i.e. 

\[ X(t, x) = x + \int_0^t a(X(s,x)) ds + \int_0^t b(X(s,x)) dB(s) \]

where \(B\) is a Brownian motion with respect to a risk-neutral measure

Pathwise differentiation with respect to \(x\) yields
    \[ Y(t, x) = \frac{\partial}{\partial x} X(t, x) = 1 + \int_0^t a'(X(s,x))Y(s, x) ds + \int_0^t b'(X(s,x))Y(s,x) dB(s) \]
These computations suggest for \(z(x) := \mathbb{E}[f(X(T, x))]\):
\[ z'(x_0) = \mathbb{E}[f'(X(T,x_0))Y(T, x_0)] \], where \((X, Y)\) satisfies a coupled pair of SDEs:
\( dX(t) &= a(X(t)) dt + b(X(t)) dB(t), \quad X(0) = x_0 \)
\( dY(t) &= a'(X(t))Y(t) dt + b'(X(t)) Y(t) dB(t),\ Y(0) = 1\)

Explain how the setting looks like that admits an application of the likelihood ratio
method. How is the score function defined? What is the formula for the derivative
that involves the score function?


Suppose we have \( Z = f(X_1, \dots, X_n) \) and \( X \) has density \(g_\theta \), where 
\( \theta \) are the parameters of \( g \).

We then have
\( \mathbb{E}_{\vartheta}[Z] = \int f(x) g_\theta(x) dx \)
\( \frac{d}{d\vartheta} \mathbb{E}_{\vartheta}[Z] = \int f(x) \frac{d}{d\vartheta} g_\theta(x) dx\)
where we have assumed the interchange of the order of differentiation and integration is again justified. Writing
\( \dot{g}_\vartheta \) for \( \frac{dg_\vartheta}{d\vartheta} \)
this is the same as
\( \int f(x) \frac{g_\vartheta}{g_\vartheta}\dot{g}_\vartheta dx\)
\( = \mathbb{E}_{\vartheta}[f(x) \frac{\dot{g}_\vartheta}{g_\vartheta}]\)

Note that \( S_{X, \vartheta} = \frac{\dot{g}_\vartheta}{g_\vartheta} = \frac{d}{d\vartheta} \log(g_\vartheta) \), which in general 
is also referred as the score function.

We could also interpret this as a change of measures (Radon-Nikodym) theorem, by defining the measures \( d\mathbb{P}_{\vartheta} = g_{\vartheta}\)
and \( d\tilde{\mathbb{P}} = \frac{g_{\vartheta}}{\log(g_\vartheta)} \), note that \( d\mathbb{P}_{\vartheta} << d\tilde{\mathbb{P}}\), 
the Radon-Nikodym derivative is then \( \frac{d\mathbb{P}_{\vartheta}}{d\tilde{\mathbb{P}}} = \log(g_\vartheta)\), which is also referred to as
\( L(\vartheta) \) in the lecture and we have the equations \(z(\vartheta) = \mathbb{E}_{\vartheta}[Z] = \tilde{\mathbb{E}}[Z L(\vartheta)\)
and \( z'(\vartheta) = \tilde{\mathbb{E}}_[ZL'] \) as we have seen before.

the generalization to higher dimensions is simply to set the score function to \( S_{X,\vartheta} = \sum_{i=1}^{d} \frac{d}{d\vartheta} \log(f^{i}(X) \)
where \( f^{i} \) is the \( i \)-th component.
It holds in particular \( \tilde{\mathbb{E}}[S_{X, \vartheta}] = 0 \) due to 
\( \int \frac{g_{\vartheta}'}{g_\vartheta} g_{\vartheta} d\tilde{\mathbb{P}} = \int g_{\vartheta}' d\tilde{\mathbb{P}} \)
and \( 1 = \int g_{\vartheta} d\tilde{\mathbb{P}} \) we can differentiate on both sides and by assumption may move the 
differential inside the integral to see the equality to the previous term.

Discuss the application of the likelihood ratio method in an insurance risk model.

Consider insurance risks modeled by a compound Poisson sum:
\( \mathcal{C} = \sum_{i=1}^{N} V_i\) where the \( (V_i)_{i \in \mathbb{N}} \) are i.i.d.
and \( N \sim \text{Poisson}(\lambda) \).

a) 
Define \( z(\lambda) = \mathbb{P}[C > c] \). Aim: Estimate \( z'(\lambda) \).

* Method: Reference measure is counting measure, 
\( f_{\lambda}(x) = e^{-\lambda} \frac{\lambda^x}{x!} \)
Then \( S_N(\lambda) = \frac{d}{d \lambda} \log f_{\lambda}(N) = \frac{N}{\lambda} - 1 \) 
hence \( z'(\lambda_0) = \mathbb{E}_{\lambda_0} [ 1_{\{C > c\}} \cdot (\frac{N}{\lambda_0} - 1) ] \)

b) 
Assume now that \(\lambda\) is fixed. Instead we suppose that the \(V_i\)s have a density \(f_{\vartheta}\) with respect to the Lebesgue measure.
Define \(z(\vartheta) = \mathbb{P}[C > c] \). Aim: Estimate \(z'(\vartheta)\)!
Method: Reference measure for individual \(V_i\) is Lebesgue measure. We define
\( S_i = \frac{d}{d \vartheta} \log f_{\vartheta}(V_i) \)
Setting \( S := \sum_{i=1}^N S_i \), one obtains \( z'(\vartheta_0) = \mathbb{E}_{\vartheta_0}[1_{\{C > c\}} \cdot S] \).

Explain the computation of vega in the Black-Scholes model using the likelihood
ratio method.

We analyze the dependence on the volatility in the Black-Scholes model.
The price of a European call option is
\( z(\sigma)=\mathbb{E}[e^{-rT}(S(0)e^{X}-K)^{+}] \text{ with } X\sim\mathcal{N}((r-\frac{\sigma^{2}}{2})T,\sigma^{2}T)\)
Aim: Estimate \( z'(\sigma_{0})! \) ('Vega')
Methods:
\( S_{X}(\sigma)=\frac{d}{d\sigma}log(\frac{1}{\sqrt{2\pi T}\sigma}exp(-\frac{(x-(r-\frac{\sigma^{2}}{2})T)^{2}}{2\sigma^{2}T}))=-\frac{1}{\sigma}+\frac{x^{2}}{\sigma^{3}T}-\frac{2xr}{\sigma^{3}}-\frac{-4\sigma^{3}(r-\frac{\sigma^{2}}{2})T-4\sigma(r-\frac{\sigma^{2}}{2})^{2}T}{4\sigma^{4}}.\)
Then:
\( z'(\sigma_{0})=\mathbb{E}[e^{-rT}(S(0)e^{X}-K)^{+}\cdot S_{X}(\sigma_{0})] \text{ with } X\sim\mathcal{N}((r-\frac{\sigma_{0}^{2}}{2})T,\sigma_{0}^{2}T) \)




\(\tilde{\pi}_{x_{-i}}^{(i)}(x_i) = \frac{\tilde{\pi}(x)}{\sum_{y_i : (y_i, x_{-i}) \in E} \tilde{\pi}(y_i, x_{-i})}\)

The Gibbs sampler relies on a specific structure of the state space. Please specify
this. How can the conditional distribution of a component given the remaining
ones be computed in this setting? State the algorithm of the Gibbs sampler.
Under suitable conditions, what is known about the invariant distribution and
convergence of the Gibbs sampler?

We assume that the state space has a specific structure:

\(E \subseteq E_1 \times E_2 \times \dots \times E_d\)

i.e., each state is characterized by its \(d\) components.

The idea of the Gibbs sampler is that only one component is updated in each step while all others remain the same. 
This leads to a particularly simple recursive structure that can easily be implemented algorithmically.

We introduce the following notation for a given \(x \in E\)

\(
x = (x_1, ..., x_d),\ x_1 \in E_1, x_2 \in E_2, \dots, x_d \in E_d 
\)
\(x_{-j} = (x_1, ..., x_{j-1}, x_{j+1}, \dots, x_d)\)

Define \( x \sim_j y \) if \( x_i = y_i \) for all \( i \neq j \) and let \( \pi(\vartheta|y) \) be the usual
posterior distribution for \( \vartheta \), given \( y \).

Then the conditional transition probability is defined as
\( \tilde{\pi}(\vartheta_{-j}|\vartheta) = \tilde{\pi}(\vartheta_{j}|\vartheta_1 = \theta_1, \dots, \vartheta_{j-1} = \theta_{j-1}, \vartheta_{i+1} = \theta_{i+1}, \dots, y) \), which coincides in this situation with the likelihood (conjugate priors)
\( p_j(\vartheta) = (d\sum_{z \in \Theta; z \sim_j x} \pi(z))^{-1}\pi(y) \) if \( x \sim_j y \) and \( 0 \) otherwise

Let \( \eta^{0} = (\vartheta_1, \dots, \vartheta_d) \) be the initial state then the Gibbs sampling proceeds by repeating the following 2 steps:
1. Pick an index \( k \in \{1, . . . , m\} \) either via round-robin or uniformly at random
2. Update \(\eta^{t}\) to \( \eta^{t+1} \) according to the distribution \(p_k(\eta)\)

It can be shown that this generates a Markov chain with the transition probabilities of \( p_j(\vartheta) \)

In the case that the indices are chosen with uniform distribution this is a special case  
of the Metropolis-Hastings algorithm with \( Q_k(y|x) = p(y|x)) \) if \( y_ \sim_j x \) and \( 0 \) otherwise.

The (prior) distribution \(\pi(x)\) is invariant for the Gibbs sampler.
If \(E\) is finite and the Gibbs sampler is irreducible, then \(d_{n}(x)\rightarrow0\), i.e.,
\(\mathcal{L}(\xi_{n}) \rightarrow \pi\)
in total variation as when \( n \to \infty \) as well as
\(\hat{f_{n}} := \frac{1}{n}\sum_{k=1}^{n} f(\xi_{k}) \xrightarrow \pi(f)\)
a.s. and in \( L^{1} \)


Consider a hidden Markov model. What is the conditional distribution of a com-
ponent given the remaining ones and given the observation? Formally derive this
result.

An application of the Gibbs sampler are hidden Markov models:
Observation: \(Y_{0}, Y_{1}, ..., Y_{n}\)
Hidden Markov process: 
\(X_{0}, X_{1}, ..., X_{n}\)
Known structure and known quantities
* Initial distribution and transition probabilities \(q(x_{i-1}, x_{i})\) of the hidden Markov chain
* Conditional independence of \((Y_{i})\) given \(X\); conditional distribution of \(Y_{i} = y_{i}\) given \((X_{i})_{i} = (x_{i})_{i}\) is \(g_{x_{j}}(y_{i})\)

These specifications admit a representation of the conditional distribution of the hidden \((X_{i})\) given the observed \((Y_{i})\) as:

\( \mathbb{P}_{\nu}(X_{0} = x_{0}, X_{1} = x_{1}, ..., X_{n} = x_{n} | Y_{0} = y_{0}, Y_{1} = y_{1}, ..., Y_{n} = y_{n}) \)

\( = \frac{\nu(x_0) q(x_0, x_1) \dots q(x_{n-1}, x_n) g_{x_0}(y_0) g_{x_1}(y_1) \dots g_{x_n}(y_n)}{\mathbb{P}(Y_0 = y_0, ..., Y_n = y_n)} \)

The aim is to sample from this conditional distribution!
For the Gibbs sampler one needs to deal with terms of the form:

\( p_i(\vartheta) = (d\sum_{z \in \Theta; z \sim_i x} \pi(z))^{-1}\pi(y) \) if \( x \sim_i y \) and \( 0 \) otherwise

One needs to distinguish \(i = 0, 0 < i < n, i = n\). If one focuses on \(0 < i < n\) most terms cancel out and the simple result is:

\(p_i(x) = \frac{q(x_{i-1}, x_{i}) q(x_{i}, x_{i+1}) g_{x_{i}}(y_{i})}{\sum_{z_{i} : (z_{i}, x_{-i}) \in E} q(x_{i-1}, z_{i}) q(z_{i}, x_{i+1}) g_{z_{i}}(y_{i})}\)

Analyzing the conditional distribution of a component given the remaining ones
in the Ising model shows that at a chosen site the Gibbs sampler flips the state
with a conditional Bernoulli probability. Derive the formula for this conditional
probability used in the sampler.

Let \( \{x_i\}_{i \in S}, \{y_i\}_{i \in S} \) be a set of sides in \( S \) with \( x_i, y_i \in \{-1, 1\} \)
Observe that if \(y \sim_i x\) in this model, then \(y_i = -x_i\).

We introduce the following notation:
\(n_i(x) = \) number of neighbors \(i'\) of \(i\) with \(x_{i'} \neq x_i\)
\(m_i(x) = \) number of neighbors \(i'\) of \(i\) with \(x_{i'} = x_i\)

Since 
\(\tilde{pi}(x) = e^{-\beta n_i(x)} \cdot d_i(x) \quad \text{with} \quad d_i(x) = \exp \left\{ -\beta \sum_{\substack{j \neq i, k \neq i, \\ \{j, k\} \in E}} 1_{\{x_j \neq x_k\}} \right\},\)
the term \(d_i(x)\) cancels out in the relevant fraction.

This provides us with the following result for \(x \sim_i y\):

\(p(y,x_{-i}) = p_i(y) = \frac{e^{-\beta m_i(x)}}{e^{-\beta n_i(x)} + e^{-\beta m_i(x)}}\)

The Gibbs sampler in the Ising model operates as follows:
1. Choose a site \(i\) (systematic scan or random scan).
2. Flipping is a Bernoulli argument with the probability computed in equation of \( p_i(y) \)
